% ===== PRÉAMBULE CANONIQUE — à coller VERBATIM en tête de CHAQUE .tex =====
\documentclass[11pt,a4paper]{article}
\usepackage[utf8]{inputenc}\usepackage[T1]{fontenc}\usepackage{lmodern}
\usepackage[french]{babel}
\usepackage{amsmath,amssymb,mathtools}\usepackage{icomma}
\usepackage{siunitx}\sisetup{locale=FR}  % \qty{}{}, \si{}, \ang{}
\usepackage{xcolor}\usepackage{colortbl}
\usepackage{tikz}\usepackage{pgfplots}\pgfplotsset{compat=1.18}
\usepackage[siunitx,europeanresistors]{circuitikz} % circuits (élec.)
\usepackage[most]{tcolorbox}\usepackage{enumitem}\usepackage{array,booktabs}
\usepackage[a4paper,margin=2.2cm]{geometry}
\usetikzlibrary{arrows.meta,calc,angles,quotes,positioning,patterns,%
  backgrounds,shapes.geometric,decorations.pathmorphing,decorations.markings,intersections}
\definecolor{primary}{RGB}{222,49,99}\definecolor{primarydark}{RGB}{150,22,55}
\definecolor{defcol}{RGB}{0,114,178}\definecolor{methcol}{RGB}{52,92,148}
\definecolor{excol}{RGB}{0,135,90}\definecolor{attncol}{RGB}{200,40,55}
\tcbset{boxbase/.style={enhanced,breakable,boxrule=0.9pt,arc=2.5pt,
  left=3mm,right=3mm,top=2.3mm,bottom=2.3mm,fonttitle=\bfseries,coltitle=white}}
% --- boîtes TUTORIEL (noms EXACTS, ne pas renommer) ---
\newtcolorbox{cle}[1][]{boxbase,colback=defcol!6,colframe=defcol,
  title={\textsc{À retenir}\ifx\relax#1\relax\relax\else\textmd{ : \itshape #1}\fi}}
\newtcolorbox{methode}[1][]{boxbase,colback=methcol!7,colframe=methcol,sharp corners,
  title={\textsc{Méthode}\ifx\relax#1\relax\relax\else\textmd{ --- \itshape #1}\fi}}
\newtcolorbox{exemplebox}[1][]{boxbase,colback=excol!5,colframe=excol,
  title={\textsc{Exemple}\ifx\relax#1\relax\relax\else\textmd{ : \itshape #1}\fi}}
\newtcolorbox{piege}[1][]{boxbase,colback=attncol!6,colframe=attncol,
  title={\textsc{Piège}\ifx\relax#1\relax\relax\else\textmd{ : \itshape #1}\fi}}
\newtcolorbox{remarque}[1][]{boxbase,colback=black!4,colframe=black!45,
  title={\textsc{Remarque}\ifx\relax#1\relax\relax\else\textmd{ : \itshape #1}\fi}}
% --- boîtes EXERCICE / CORRIGÉ (noms EXACTS, auto-comptées) ---
\newtcolorbox[auto counter]{exo}[1][]{boxbase,colback=primary!4,colframe=primary,
  title={\textsc{Exercice}~\thetcbcounter\ifx\relax#1\relax\relax\else\textmd{ --- \itshape #1}\fi}}
\newtcolorbox[auto counter]{corrige}[1][]{boxbase,colback=excol!4,colframe=excol,
  title={\textsc{Corrigé}~\thetcbcounter\ifx\relax#1\relax\relax\else\textmd{ --- \itshape #1}\fi}}
\newcommand{\vv}[1]{\vec{#1}}\newcommand{\ds}{\displaystyle}
\newcommand{\R}{\mathbb{R}}\newcommand{\N}{\mathbb{N}}\newcommand{\Z}{\mathbb{Z}}
% (un \usepackage{fancyhdr}+en-tête est possible ; il est ignoré côté web)
% =========================================================================
\begin{document}
\sisetup{per-mode=symbol}

\begin{corrige}[vrai ou faux : la voiture et le camion]
Le camion ($\qty{4}{\tonne}=\qty{4000}{\kilogram}$) est deux fois plus lourd que la voiture, à la même
vitesse : $p_{\text{camion}}=2\,p_{\text{voiture}}$.
\begin{enumerate}[label=\alph*)]
  \item \textbf{Faux}. À vitesse égale, $p=mv$ est proportionnelle à la masse : le camion a une
        quantité de mouvement \emph{deux fois} plus grande.
  \item \textbf{Vrai}. Le camion a la plus grande quantité de mouvement : il est le plus difficile à
        arrêter (force ou distance de freinage plus grande).
  \item \textbf{Faux}. C'est le contraire : la voiture, moins « lancée », est plus facile à arrêter.
  \item \textbf{Vrai}. $p=mv$ augmente avec $v$ : plus vite $\Rightarrow$ plus grand $p$ $\Rightarrow$
        plus difficile à arrêter.
\end{enumerate}
\end{corrige}

\begin{corrige}[même énergie cinétique, quelle quantité de mouvement ?]
\begin{enumerate}[label=\alph*)]
  \item $v_1=\sqrt{\dfrac{2E_c}{m_1}}=\sqrt{\dfrac{2\times64}{2}}=\sqrt{64}=\qty{8}{\metre\per\second}$,
        d'où $p_1=m_1 v_1=2\times8=\qty{16}{\kilogram\metre\per\second}$.
  \item $v_2=\sqrt{\dfrac{2\times64}{8}}=\sqrt{16}=\qty{4}{\metre\per\second}$, d'où
        $p_2=m_2 v_2=8\times4=\qty{32}{\kilogram\metre\per\second}$. À énergie cinétique égale, c'est le
        \emph{plus massif} qui a la plus grande quantité de mouvement ($p_2=2\,p_1$).
\end{enumerate}
\end{corrige}

\begin{corrige}[relier énergie cinétique et quantité de mouvement]
\begin{enumerate}[label=\alph*)]
  \item De $p=mv$ on tire $v=\dfrac{p}{m}$. En reportant dans $E_c=\tfrac12 m v^2$ :
        \[
        E_c=\tfrac12 m\left(\frac{p}{m}\right)^2=\tfrac12 m\,\frac{p^2}{m^2}=\frac{p^2}{2m}.
        \]
  \item $E_c=\dfrac{p^2}{2m}=\dfrac{20^2}{2\times5}=\dfrac{400}{10}=\qty{40}{\joule}$.
        Vérification : $v=\dfrac{p}{m}=\dfrac{20}{5}=\qty{4}{\metre\per\second}$, donc
        $E_c=\tfrac12\times5\times4^2=\qty{40}{\joule}$ — cohérent.
\end{enumerate}
\end{corrige}

\begin{corrige}[vrai ou faux : propriétés de la quantité de mouvement]
\begin{enumerate}[label=\alph*)]
  \item \textbf{Vrai}. $\vv p=m\vv v$ : c'est un vecteur (direction et sens de $\vv v$).
  \item \textbf{Faux}. Son unité est le \si{\kilogram\metre\per\second} (le joule est l'unité d'énergie).
  \item \textbf{Vrai}. À $v$ égal, $p=mv$ croît avec la masse : le plus massif a le plus grand $p$.
  \item \textbf{Faux}. À énergie cinétique égale, $p=\sqrt{2mE_c}$ dépend de la masse : deux objets de
        masses différentes ont des quantités de mouvement différentes (cf. exercice précédent).
\end{enumerate}
\end{corrige}
\end{document}
